\documentclass[11pt]{article}
\usepackage[a4paper,left=25mm,right=25mm,top=25mm,bottom=25mm]{geometry}
\usepackage[utf8]{inputenc}
\usepackage[square]{natbib}
\usepackage{fancyhdr}
\usepackage[pdftex]{graphicx}
\usepackage{subcaption}
\captionsetup{font=sf, labelfont={sf,bf}}
\renewcommand{\figurename}{Şekil}
\renewcommand{\tablename}{Tablo}
%\pdfcompresslevel=9
%\def\shell_escape{1}

\def\bildirino{UHUK-2026-000}
\def\Baslik{YAPAY ZEKÂ TEMELLİ KESTİRİMCİ BAKIM: "IDG" FREKANS TAHMİNİ UYGULAMASI}
\def\yazarlar{KIZILDOĞAN, THYYAZAR1, THYYAZAR2 ve THYYAZAR3}  %..Sadece yazarların soyadları
\author{                                 %..Yazarlar
    \parbox[t]{5cm}{\lsf\centering
            Göksu Anıl Kızıldoğan\thanks{\ssf Analythinx,
                            E-posta: goksu.kizildogan@analythinx.com}\\
            \sf
            Analythinx, Eskişehir}
    \hspace*{4mm}
    \parbox[t]{5cm}{\lsf\centering
            ThyYazar1\thanks{\ssf Turkish Technology,
                            E-posta: thyyazar1@thy.com}
            ~,
            Yazar 2\thanks{\ssf Turkish Technology,
                            E-posta: thyyazar2@thy.com}
            ~ve
            Yazar 3\thanks{\ssf Turkish Technology,
                            E-posta: thyyazar3@thy.com}\\
            \sf
            Turkish Technology, İstanbul}
}

%%%%%...Tanımlamalar...%%%%%%%%%%%%%%%%%%%%%%%%%%%%%%%%%%%%%%%%%%%%%%%%%%%%%
\font\sf=cmss10  scaled 1100
\font\ssf=cmss10 scaled  900
\font\lsf=cmss10 scaled 1200
\font\bsf=cmssbx10 scaled 1100
\font\lbsf=cmssbx10 scaled 1200
\parindent 0pt \partopsep 3mm plus 1pt minus 1pt \parskip 3mm plus 1pt 
\headheight 16pt
\def\baslik#1{\vspace{2\abovedisplayskip}\centerline{\bsf #1}\nobreak\sloppy}
\def\altbaslik#1{\vspace{\partopsep}{\bsf #1}\par\nobreak\sloppy}
\def\altaltbaslik#1{\par\sf\underline{#1: }\nobreak}
\def\refname{Kaynaklar}
\lhead{\ssf \yazarlar}
\rhead{\ssf \bildirino}
\cfoot{\sf \thepage \\ \ssf Ulusal Havacılık ve Uzay Konferansı}
\pagestyle{fancy}
\begin{document}
\thispagestyle{empty}
\date{}
\title{\sf 11. ULUSAL HAVACILIK VE UZAY KONFERANSI \hfill 
       \lbsf \bildirino\\[-2mm]
       \sf 17-19 Eylül 2026, Eskişehir Osmangazi Üniversitesi, Eskişehir \hfill~ \\[16mm]
       \lbsf \Baslik }
\maketitle
\raggedright
\sf
%%%%%%%%%%%%%%%%%%%%%%%%%%%%%%%%%%%%%%%%%%%%%%%%%%%%%%%%%%%%%%%%%%%%%%%%%%%%%


\baslik{ÖZET}

{\it Tek parağraflık özet bölümü italik fontla
yazılmalı, 250 sözcüğü geçmemelidir. Bildiride
sunulan çalışmanın amacı, kullanılan yöntem ve 
elde edilen sonuçlar kısaca belirtilmelidir. }

\baslik{GİRİŞ}

Bildiri özetleri ve bildiriler A4 boyutundaki
sayfalara, tek sütun, tek satır aralıklı olarak,
11 punto, Sans Serif fontlarla (Helvetica, Arial,
Tahoma, Verdana) hazırlanmalıdır. Tüm kenar boşlukları
25mm, metin sola dayalı (left justified) olarak hece
bölmeden yazılmalıdır. Paragraflar arası 3mm boşluk
bırakılmalıdır. Sayfa sayısında bir sınırlama yoktur.

Ana başlık ilk sayfanın üst kısmında ortalanarak 12
punto, kalın Sans Serif fontla yazılmalıdır. Yazarların
isim/kurum/şehir bilgileri her yazar için 3-satırlık
bloklar halinde ana başlığın altında yatay olarak yer
almalıdır. Yazarların ünvanları ve e-posta adresleri
dipnotta verilmelidir. Kabul edilen bildirilere verilen
Bildiri Numaraları ana başlığın üzerine kalın fontla
eklenmelidir.

Bölüm Başlıkları sayfada ortalanmalı, ve kalın punto,
büyük harflerle yazılmalıdır.

Üstbilgi satırına Bildiri Numarası (kabul edilen
bildiri özetlerine verilen) ve yazarların soyadları
yazılmalıdır. Altbilgi satırında Ulusal Havacılık ve Uzay
Konferansı yazısı yer almalıdır.

Şekil ve çizelgeler ortalanmalı (Şekil \ref{fig1}),
ve metnin içinde verilmelidir.

\begin{figure}[t]\centering
%\includegraphics{sekil.pdf} \\
\caption{Örnek şekil}
\label{fig1}
\end{figure}
\newpage

\altbaslik{İkincil Başlık}

Birinci altbaşlıklar sola dayalı olarak kalın punto ile ve
kelimelerin ilk harfi büyük olarak yazılmalıdır.

\altaltbaslik{İkinci  Altbaşlık} İkinci  altbaşlıklar sola dayalı
olarak ve altı çizili olarak verilmeli, metin hiç ara
vermeden, başlığın bittiği yerden başlayarak yazılmalıdır.


\baslik{YÖNTEM}

Kaynaklar yazar-tarih yaklaşımı ile verilmelidir\citep{ref1}.
Bildiri metninde kaynaklar köşeli parantez içerisinde 
\citep{ref1,ref2}, Kaynaklar listesindeki alfabetik olarak yazar-tarih 
sıralaması ile gösterilir.

\baslik{UYGULAMALAR VE DEĞERLENDİRME}

Çalışmada yapılan tüm uygulamalar, elde edilen detaylı
sonuçlar, karşılaştırmalar, ve ilgili yorumlar
bu bölümde yer almalıdır.


\baslik{SONUÇ}

Bu bölümde çalışmanın kısa bir değerlendirmesi yapılmalı,
ulaşılan temel sonuçlara ve ileriye dönük tavsiyelere
yer verilmelidir.

%\newpage
\begin{thebibliography}{99}
\bibitem[SoyisimA(2022)]{ref1}
   SoyisimA, İ., 2022.
   {\it Yazar soyisimlerine göre alfabetik sırada  ilk referans..},
   Konferans, Yer, 12-14 Haziran
\bibitem[SoyisimB, SoyisimC ve SoyisimD(2021)]{ref2}
   SoyisimB, J., SoyisimC, K. ve SoyisimD. L.,  2021.
   {\it Yazar soyisimlerine göre alfabetik sırada  ikinci referans..},
   Yayımcı Kuruluş,  Cilt.2, s.100-119

\end{thebibliography}

\end{document}
